%! Solving Radical Equations \& Extraneous Solutions — Unit 5, Lesson 5.5 — Student Packet Cover Sheet
\documentclass[10pt]{article}
\usepackage{algebra2-article}
\usepackage{algebra2-boxes}
\usepackage{ltablex}
\keepXColumns

\begin{document}

% Full-bleed forest banner
\begin{tikzpicture}[remember picture, overlay]
  \fill[forest] (current page.north west)
    rectangle ([yshift=-0.9in]current page.north east);
\end{tikzpicture}
\vspace{-0.8in}

\begin{center}
  {\color{white}\textbf{\LARGE Algebra 2}}\\[4pt]
  {\color{forestmid}\large\textbf{Unit 5 \quad Radical Functions}}\\[6pt]
  {\color{white}\textbf{\large Lesson 5.5 \quad Solving Radical Equations \& Extraneous Solutions}}
\end{center}

\vspace{0.22in}
\namedateperiod
\vspace{0.18in}

\begin{learningtargetbox}
\small
\begin{enumerate}[leftmargin=1.5em, itemsep=5pt]
  \item \ldots{} solve an equation containing one radical by \textbf{isolating} the radical and raising both sides to the power equal to the index --- and solve an equation of the form $x^{m/n}=k$ by raising both sides to the reciprocal power.
  \item \ldots{} \textbf{check every candidate} in the original equation, name the \textbf{extraneous} ones, and recognize a ``no real solution'' equation before doing any algebra at all.
  \item \ldots{} \textbf{verify} a solution set on a graph, by counting where the two sides of the equation intersect.
  \item \ldots{} \textbf{justify} why an extraneous solution appears for an even index --- and why it can never appear from the raising step when the index is odd.
\end{enumerate}
\end{learningtargetbox}

\vspace{0.14in}

\begin{tocbox}
\small
\renewcommand{\arraystretch}{1.5}
\begin{tabularx}{\linewidth}{@{} c l X r @{}}
\rowcolor{forestbg}
\textbf{\#} & \textbf{Component} & \textbf{Description} & \textbf{Score} \\
\midrule
1 & Warm-Up        & Spiral review: undoing a radical, the one-way street of squaring, and one quadratic you will meet again & \blank{1.2cm} \\
2 & Guided Notes   & The four-step method, why extraneous solutions exist, the sign shortcut, odd indices, graphs, and $x^{m/n}=k$ & \blank{1.2cm} \\
3 & Group Activity & \emph{Candidate or Solution?} --- solving, checking, and sorting, by tier & \blank{1.2cm} \\
4 & Exit Ticket    & Quick check: solve, sort the candidates, plus an SOL-style item & \blank{1.2cm} \\
5 & Homework       & Mixed equations, rational exponents, a graph, an error to diagnose, and reading the skid marks & \blank{1.2cm} \\
\midrule
  & \multicolumn{2}{r}{\textbf{Total}} & \blank{1.2cm} \\
\end{tabularx}
\end{tocbox}

\vspace{0.10in}

\begin{remindbox}
\small For four lessons you have been \emph{building} radical expressions and graphs. Today you take
one apart. To undo an $n$th root, raise both sides to the $n$th power --- but only after you
\textbf{isolate} the radical, because squaring erases a radical only when it is standing alone
($\left(\sqrt{x}+3\right)^{2}$ is still full of radicals). Then comes the step that makes this
lesson different from every other equation-solving lesson you have had: \textbf{you must check}.
Squaring is a one-way street. If $a=b$ then $a^{2}=b^{2}$ always, but $a^{2}=b^{2}$ only tells you
that $a=b$ \emph{or} $a=-b$ --- so squaring can hand back answers the original equation never
accepted. Those are \textbf{extraneous solutions}, and they are not random: a candidate fails
exactly when it makes the \emph{other} side of the equation negative, because a principal square
root never is. That also gives you a free shortcut --- an isolated even root set equal to a negative
number has \emph{no real solution}, and you can say so before writing a single line. Change the
index and the whole problem changes, as it has all unit: cubing keeps the sign of what it acts on,
so a cube root equation can never manufacture a phantom. Finally, every one of these questions is a
question about Lesson 5.4's graphs --- solving $\sqrt{x+4}-1=3$ is asking where that curve crosses
the line $y=3$, and the number of crossings is the number of real solutions.
\end{remindbox}

\end{document}
